\begin{figure}[ht]
\centering
\begin{tikzpicture}[x=0.6cm, y=0.45cm, font=\footnotesize]

% Log-log: x = log10(h), y = log10(|error|)
% h from 10^0 down to 10^{-15}
% Truncation error decreases as h^p; round-off increases as 1/h for
% finite-difference derivative; total has a U shape with minimum near
% h ~ sqrt(u) ~ 10^{-8} for first-derivative central difference.

% Axes
\draw[->, thick] (-15.5,0) -- (1,0) node[right] {$\log_{10} h$};
\draw[->, thick] (-15.5,-16) -- (-15.5,1) node[above, align=center]
  {$\log_{10}\abs{\text{error}}$};

% Tick marks
\foreach \k in {0,-3,-6,-9,-12,-15} {
  \draw (\k,0) -- (\k,0.3);
  \node[anchor=south, font=\scriptsize] at (\k,0.3) {$10^{\k}$};
}
\foreach \d in {0,-4,-8,-12,-16} {
  \draw (-15.5,\d) -- (-15.8,\d);
  \node[anchor=east, font=\scriptsize] at (-15.8,\d) {$10^{\d}$};
}

% Truncation error: O(h^2) for central difference, slope +2
% y = 2*x - 2 (so at h=10^0 error = 10^{-2}; at h = 10^{-8} error = 10^{-18})
\draw[thick, engblue] (0,-2) -- (-8,-18);
\node[engblue, font=\scriptsize, anchor=west] at (-3,-7)
  {truncation $\sim h^{2}$};

% Round-off: ~ u / h, slope -1 (here y = -x - 16; at h=10^0 error = 10^{-16}, at h=10^{-15} error = 10^{-1})
\draw[thick, engred] (0,-16) -- (-15,-1);
\node[engred, font=\scriptsize, anchor=east] at (-12,-4)
  {round-off $\sim u/h$};

% Total error: pointwise max of the two
\draw[thick, black, dashed] plot coordinates {
  (0,-2) (-1,-4) (-2,-6) (-3,-8) (-4,-10) (-5,-12) (-5.5,-12.5)
  (-5.8,-12.7) (-6,-12.4) (-7,-11.8) (-8,-10.7) (-10,-8.5)
  (-12,-6.5) (-15,-1)
};
\node[font=\scriptsize, anchor=south] at (-6,-12.5)
  {optimal $h^{*} \approx u^{1/3}$};

% Mark optimum (h ~ 10^{-5.3}, error ~ 10^{-10.6}) for central diff
\fill[black] (-5.5,-12.5) circle (2pt);

\end{tikzpicture}
\caption[Truncation vs round-off in finite-difference differentiation]{The
total error of a finite-difference approximation to a derivative as a
function of step size $h$, on log-log axes. The truncation error of the
centred difference $f'(x) \approx (f(x + h) - f(x - h))/(2 h)$ scales
as $h^{2}$ and decreases with $h$. The round-off error of the
subtraction in the numerator scales as $u / h$ and increases as $h$
shrinks. The total error has a minimum near $h^{*} \approx u^{1/3}$,
with $u$ the unit round-off; for binary64, $h^{*} \approx 10^{-5.3}$,
producing an attainable error near $10^{-10.6}$. A naive engineer who
chases truncation error down past this point by shrinking $h$ further
loses digits to round-off and reports a worse derivative than at the
optimum. Source: authors' analysis; the $u^{1/3}$ optimum is standard,
see \cite[\S 5.7]{text:higham-numerical}.}
\label{fig:vol02:ch16:error-vs-h}
\end{figure}
